% TOP_PROBLEM: {"id": 2304003, "title": "Research Problems in Function Theory — Problem 4.3", "category_id": 8, "category": {"id": 8, "name": "analysis", "display_name": "Analysis", "description": "Limits, continuity, calculus, and function theory.", "slug": "analysis", "order_index": 8, "created_at": "2026-07-31T15:26:25.671Z"}, "status": "open"}
% FIRST_POSTED: null
% ENTRY_KIND: statement_only
% Imported from: batch13.tex; UnsolvedMath ID 2304003
\documentclass[11pt]{book}
\usepackage{fontspec}
\setmainfont{Latin Modern Roman}
\setsansfont{Latin Modern Sans}
\usepackage{amsmath,amssymb,mathtools}
\usepackage{unicode-math}
\setmathfont{Latin Modern Math}
\usepackage{xurl}
\usepackage[hidelinks]{hyperref}
\newcommand{\sref}[2]{\hyperlink{source:#1:#2}{[S#2]}}
\newcommand{\cataloguescope}[1]{\paragraph{Mathematical question.}#1}
\begin{document}
% UnsolvedMath ID 2304003; source catalogue code AMR-022-4003
\setcounter{section}{2304002}
\section{Research Problems in Function Theory — Problem 4.3}\label{2304003}
{\small\sffamily UnsolvedMath ID 2304003 · Analysis}
\subsection{Definitions and mathematical statement}

\paragraph{Shared notation.}$\mathbb N=\{1,2,\ldots\}$, and $\mathbb Z,\mathbb Q,\mathbb R,\mathbb C$ denote the integers, rationals, reals and complex numbers. Any other conventions are specified locally.


Write $\mathbb D=\{z\in\mathbb C:|z|<1\}$ and $\mathbb T=\partial\mathbb D$. For $N\ge1$, let
\[P_N(z)=\sum_{j=1}^Na_jz^{m_j},\qquad 0\le m_1<\cdots<m_N,\quad a_j\in\mathbb C\setminus\{0\},
\]
be a polynomial with exactly $N$ nonzero terms, normalized by $\max_{|z|=1}|P_N(z)|\le1$. For $1\le n\le N$, its initial partial sum is $P_n(z)=\sum_{j=1}^na_jz^{m_j}$. Denote by $C(N,n)$ the supremum of $\max_{|z|=1}|P_n(z)|$ over all such polynomials and exponent lists.
\paragraph{Statement recorded in the supplied source.}
Determine the sharp bounds $C(N,n)$, and their growth as the number of terms increases. The exponents need not be consecutive: the source fixes a number of terms, rather than imposing a degree bound. \sref{2304003}{2}
\subsection{Short English statement}
How large can an initial partial sum of an N-term polynomial be on the unit circle when the full polynomial is bounded there by one?
\subsection{Sources}
\begin{description}
\item[\hypertarget{source:2304003:1}{S1}] ULAM.AI and the UnsolvedMath Contributors, \emph{UnsolvedMath: A Curated Collection of Open Mathematics Problems}, record 2304003 (AMR-022-4003). CC BY 4.0. \url{https://huggingface.co/datasets/ulamai/UnsolvedMath}.
\item[\hypertarget{source:2304003:2}{S2}] Walter K. Hayman and Edward F. Lingham, \emph{Research Problems in Function Theory}, arXiv:1809.07200, 2018. Problem 4.3, its complete Update 4.3, and the relevant chapter notation. \url{https://arxiv.org/pdf/1809.07200}.
\end{description}
\label{end:2304003}
\end{document}
